\documentclass[11pt]{article}
\usepackage[margin=0.72in]{geometry}
\usepackage{booktabs,longtable,hyperref}
\title{GEO Monitoring Consulting Report}
\begin{document}
\maketitle
\section*{佳农股份}
\textbf{Generated:} 2026-05-07 11:16:28 UTC\\
\textbf{Total DeepSeek calls:} 62\\
\textbf{Monitoring probes:} 60\\
\textbf{Brand mention rate:} 38.33\\
\textbf{Brand recommendation rate:} 1.67\\
\section*{Executive Summary}
佳农股份在GEO监控中品牌可见性中等，但认知度和推荐率较低，面临美菜网、快驴进货等竞品的显著压力。信息透明度不足，存在多处事实风险，需加强官方信息发布和信任信号建设。
\section*{Dimension Scorecard}
\begin{tabular}{lr}\toprule Dimension & Score\\\midrule
Visibility & 21.6\\
Inclusion & 16.7\\
Cognition & 15.9\\
Outcome & 8.8\\
\bottomrule\end{tabular}
\section*{Top Findings}
\begin{itemize}
\item 品牌提及率38.33\%，但推荐率仅1.67\%，说明品牌被提及但未被积极推荐。
\item 竞品美菜网和快驴进货在提及和推荐上均领先，佳农股份在竞争中处于劣势。
\item 36.67\%的回答需要第一方来源，表明官方信息覆盖不足。
\item 13个事实风险标志，主要集中在联系方式、价格、配送范围等关键信息上。
\item 认知维度得分15.92，低于可见性和包容性，说明用户对品牌了解不深。
\end{itemize}
\section*{90-Day Action Roadmap}
\begin{itemize}
\item 加强官方信息发布，特别是官网、联系方式、价格表等，减少对第三方来源的依赖。
\item 提升品牌推荐率，通过客户案例、品质认证、合作案例等增强信任。
\item 针对竞品优势领域（如配送时效、服务范围）进行差异化宣传。
\item 监控并回应事实风险，确保公开信息准确及时。
\item 增加在行业报告、媒体中的曝光，提升认知度。
\end{itemize}
\section*{Methodology Note}
基于60次DeepSeek调用，覆盖品牌、竞品、行业等8类查询，采样逻辑确保各维度代表性。数据聚合后计算维度得分和证据指标，风险标志基于回答中的事实不确定性。
\end{document}